\chapter*{Vorwort … zu den Lerneinheiten »Wirtschaftsmathematik«}
\addcontentsline{toc}{chapter}{Vorwort … zu den Lerneinheiten »Wirtschaftsmathematik«}

Niemand bezweifelt, dass bei der Bewältigung des täglichen Lebens die Mathematik eine große Rolle spielt. Das beginnt schon beim Preisvergleich im Supermarkt, und beim Vergleich von Hypothekenkonditionen muss man bereits etwas tiefer in die »Mathematik-Kiste« greifen. Umso mehr gilt dies, wenn man mit anspruchsvolleren Problemen aus der Berufswelt oder der Wissenschaft zu tun hat. Weder Marktforschung noch die Lösung von Optimierungsproblemen in Lagerhaltung und Logistik und schon gar nicht die Simulation von Crashtests auf Computern sind ohne mathematische Modellierung denkbar. Auch das Design neuer Werkstoffe in den Materialwissenschaften oder die Vorhersage von bevorstehenden Klimaänderungen aufgrund der Freisetzung und Ausbreitung von Treibhausgasen in der Atmosphäre basieren auf mathematisch formulierten Modellen. Die Mathematik stellt somit ein absolut notwendiges und unersetzliches Hilfsmittel der Technik und der Wirtschaftswissenschaften dar.

In der modernen Mathematik stehen statt Einzelbetrachtungen, die oft den Blick auf allgemeine Zusammenhänge verstellen, systematische Methoden im Vordergrund, die aufgrund ihrer Allgemeinheit nicht auf spezielle Anwendungsbereiche beschränkt sind. Die Konsequenz ist allerdings ein hoher Abstraktionsgrad. Um die damit verbundenen Schwierigkeiten zu vermeiden, wird in den vorliegenden Lerneinheiten auf Kosten der Exaktheit das Hauptgewicht auf Anschaulichkeit und Anwendungen gelegt. Denn für den in einem Unternehmen tätigen Ingenieur oder Betriebswirt stellt die Beschäftigung mit Mathematik keinen Selbstzweck dar. Zwar sind die Eleganz der mathematischen Darstellung und die Stringenz der Beweisführung grundlegende Kriterien der reinen Mathematik. Die im Sinne der reinen Mathematik exakte Begründung soll jedoch den fortgeschrittenen Standardlehrbüchern der Mathematik vorbehalten bleiben. Trotzdem soll in den Lerneinheiten Wert auf exakte Darstellung gelegt werden, denn exaktes Denken ist in allen Wissenschaftsbereichen wichtig. Durch konkrete Beispiele soll der Werkzeugcharakter der Mathematik und damit der Anwendungsbezug hergestellt werden. Adressat ist dabei der mathematische Anfänger, der neben schulischen Vorkenntnissen ein entwicklungsfähiges Interesse einbringen sollte.

Ziel der Lerneinheiten ist es, in einer leicht lesbaren Abhandlung die Grundlagen der Wirtschaftsmathematik darzustellen. Dabei sollen die grundlegenden mathematischen Methoden und Verfahren vorgestellt und deren Anwendung zur Lösung praktischer Aufgabenstellungen aufgezeigt werden. Die Lerneinheiten zum Modul »Wirtschaftsstatistik« werden dann der Statistik gewidmet sein.

Die Lerneinheiten dienen dem Selbststudium. Da das Lesen des Textes nicht identisch mit dem Verstehen des Stoffes ist, kommt den Beispielen und Übungsaufgaben eine besondere Bedeutung zu. Sie sollen Gelegenheit zur Wiederholung des Stoffes und zur Überprüfung des Wissens geben und dienen der Selbstkontrolle des Gelernten. Die vorliegende erste Lerneinheit enthält besonders viele Übungsaufgaben.

Die Beispiele und Übungsaufgaben sollen aber auch den Zusammenhang zwischen mathematischen Konzepten und konkreten Problemstellungen der Praxis verdeutlichen. Bei all dem gilt: Mathematik lernt man durch Tun. Die angebotenen Präsenzübungen, die zur Vertiefung des Verständnisses beitragen sollen, bieten schließlich auch die Möglichkeit zur gemeinsamen Diskussion und Lösung offener Probleme.

Diesem Ziel soll auch der abschließende Hinweis auf eine kleine Auswahl von Fachliteratur dienen:
\begin{itemize}
  \item Knorrenschild, M.: Vorkurs Mathematik. Carl Hanser Verlag, 4. Auflage 2013
  \item Poguntke, W.: Keine Angst vor Mathe. Teubner, 4. Auflage 2010
  \item Schwarze, J.: Mathematik für Wirtschaftswissenschaftler Band 1, NWB-Verlag, 14. Auflage 2015
  \item Salomon, E., Poguntke, W.: Wirtschaftsmathematik. 2. überarbeitete Auflage, Fortis, 2003
\end{itemize}

\chapter*{…zur Lerneinheit 1}
\addcontentsline{toc}{chapter}{…zur Lerneinheit 1}

Die vorliegende erste Lerneinheit behandelt Grundlagen, die für das weitere Verständnis der nächsten Lerneinheiten absolut notwendig sind. Viele von Ihnen werden große Teile des Stoffes schon einmal (in der Schule oder anderswo) kennengelernt, aber möglicherweise wieder vergessen haben. Frischen Sie Ihre Kenntnisse auf! Bearbeiten Sie die zahlreichen Übungsaufgaben! Nehmen Sie die zusätzlich angebotenen Brückenkurse und Tutorien wahr, um eventuelle Lücken zu schließen! Schauen Sie in andere Literatur! Und schließlich: Zögern Sie nicht, in den Präsenzveranstaltungen auch die »dümmsten« Fragen zu stellen! Sie wissen ja: Es gibt keine dummen Fragen …

Die Lerneinheit kann auch begleitend zu einem Brückenkurs vor oder zu Beginn des ersten Semesters verwendet werden -- bei dieser Art des Einsatzes bleibt mehr Zeit, die danach folgenden Lerneinheiten während des Semesters zu behandeln.
